\documentclass[tikz,border=5pt,12pt]{standalone}
\usepackage{resistencia}
\begin{document}
% σn(θ) y τ(θ) de una barra con σ0 = 100 MPa; datos calculados por Julia (barra.dat).
\begin{tikzpicture}
\begin{axis}[width=12.5cm,height=7.8cm,xlabel={$\theta$ (grados): ángulo de la normal con el eje},
  ylabel={tensión (MPa)},xmin=0,xmax=90,ymin=0,ymax=112,xtick={0,15,30,45,60,75,90},
  ytick={0,25,50,75,100},grid=major,grid style={gris!50},legend pos=north east,legend cell align=left]
\addplot[acero,no marks] table[x=theta,y=sn] {datos/barra.dat};
\addlegendentry{$\sigma_n=\sigma_0\cos^2\theta$}
\addplot[cobre,no marks] table[x=theta,y=tau] {datos/barra.dat};
\addlegendentry{$\tau=\sigma_0\sin\theta\cos\theta$}
\addplot[only marks,mark=*,mark size=3pt,violeta,every axis plot/.append style={},
  nodes near coords,point meta=explicit,
  nodes near coords style={font=\normalsize,text=violeta,anchor=south west,fill=white,inner sep=1pt}]
  table[x=theta,y=sn,meta=sn] {datos/barra_puntos.dat};
\addplot[only marks,mark=*,mark size=3pt,violeta,
  nodes near coords,point meta=explicit,
  nodes near coords style={font=\normalsize,text=violeta,anchor=south west,fill=white,inner sep=1pt}]
  table[x=theta,y=tau,meta=tau] {datos/barra_puntos.dat};
\end{axis}
\end{tikzpicture}
\end{document}
